\section{Farben}\label{Sec:Farben}

\subsection{Adaptive Farbe}

Die Farbdefinition \texttt{Accent} und ihre Verwandtschaft passt sich der Dokumentenoption \verb+department+\docoption{department} an und stellt so eine konsistente Farbdarstellung über das Dokument sicher - auf beim Wechsel der \texttt{department}"=Einstellung.

Hier sind die Farbnamen zur Verwendung der adaptiver Farben am Beispiel der gewählten Stiloption \verb+department=FakSG+ dargestellt.

\verb|\textcolor{Accent}{...}|\\
\textcolor{Accent}{\bfseries Die Akzentfarbe \texttt{Accent} wird automatisch an die Fakultät angepasst.}

\verb|\textcolor{Accent90}{...}|\\
\textcolor{Accent90}{\bfseries Die Akzentfarbe \texttt{Accent90} wird automatisch an die Fakultät angepasst.}

\verb|\textcolor{Accent50}{...}|\\
\textcolor{Accent50}{\bfseries Die Akzentfarbe \texttt{Accent50} wird automatisch an die Fakultät angepasst.}

\verb|\textcolor{Accent25}{...}|\\
\textcolor{Accent25}{\bfseries Die Akzentfarbe \texttt{Accent25} wird automatisch an die Fakultät angepasst.}

\verb|\textcolor{AccentFontColor}{...}|\\
\textcolor{AccentFontColor}{Die Akzentfarbe \texttt{AccentFontColor} wird automatisch an die Fakultät angepasst, wobei bei im unterschied zur Farbe \texttt{Accent} bei Farben mit mangelndem Kontrast vor weißem Hintergrund  auf \texttt{OTHgrau} ausgewichen wird.}

\begin{center}
\colorbox{Accent}{\begin{minipage}{.9\linewidth}\color{AccentContrastColor} Die Farbdefinition \texttt{AccentContrastColor} kann verwendet werden, um sicher lesbaren Text auf einen Hintergrund in Akzentfarbe zu setzen.\end{minipage}}
\end{center}

Die Dokumentenoption \verb+blackFont+ \docoption{blackFont} setzt die \texttt{AccentFontColor} ungeachtet der verwendeten Stiloption \verb+department+ auf schwarz und die \texttt{AccentContrastColor} auf weiß.
Damit lässt sich bei Nutzung der farbigen Textgestaltung z.B. in Präsentationen leicht auf ein evtl. druckfreundlicheres Farbschema wechseln.

\Wichtig{\sloppy Es wird dringend empfohlen, die vorhandenen adaptiven Farben \texttt{Accent}, \texttt{Accent90}, \texttt{Accent50}, \texttt{Accent25}, \texttt{AccentFontColor} und bei Bedarf auch \texttt{AccentContrastColor} zu verwenden.
Dadurch passt sich die Gestaltung des Dokuments automatisch korrekt an die Dokumentenoption \texttt{department} an.}


Derzeit unterstützte Werte für die Option \verb+department+\docoption{department} und ihre Bedeutung:

\begin{table}[!h!]
\centering\scriptsize
\begin{tabular}{lcccc}\toprule
\textbf{Wert} & Einrichtung & Farbschema & \verb+AccentFontColor+ & \verb+AccentContrastColor+\\\midrule
OTHR & alle & Schwarz/grau & \verb+OTHRgrau+ & \verb+black+ \\
FakA & Fakultät A & Fakultätsfarbe & \verb+OTHRgrau+ & \verb+black+ \\
%\textcolor{black!20}{FakAM} & \textcolor{black!20}{Fakultät AM} & \textcolor{black!20}{Fakultätsfarbe} & \textcolor{black!20}{Fakultätsfarbe} & \color{black!20}\verb+white+ \\
FakANK & Fakultät ANK & Fakultätsfarbe & Fakultätsfarbe & \verb+white+ \\
FakB & Fakultät B & Fakultätsfarbe & Fakultätsfarbe & \verb+white+ \\
FakBM & Fakultät BM & Fakultätsfarbe & Fakultätsfarbe & \verb+white+ \\
FakEI & Fakultät EI & Fakultätsfarbe & Fakultätsfarbe & \verb+white+ \\
FakIM & Fakultät IM & Fakultätsfarbe & Fakultätsfarbe & \verb+white+ \\
FakM & Fakultät M & Fakultätsfarbe & Fakultätsfarbe & \verb+white+ \\
FakSG & Fakultät SG & Fakultätsfarbe & Fakultätsfarbe & \verb+white+ \\
ZWW & ZWW & ZWW-Farbe & \verb+OTHgrau+ & \verb+black+ \\
Labornamen & Labore & diverse & diverse & diverse \\
\bottomrule
\end{tabular}
\end{table}

Hierbei ist die Farbe \verb+AccentFontColor+ jeweils für Text auf weißem Grund gedacht und wird im Allgemeinen auf die Farbe des mit der Stiloption  \verb+department+ gewählten Bereichs gesetzt.
Da bei manchen der vordefinierten Farben aber der Kontrast auf weißem Hintergrund zu schwach ausfällt, wird in diesen Fällen Hochschulgrau für \verb+AccentFontColor+ gewählt.

Zudem zeigt sich, dass bei Verwendung der Bereichsfarben als Hintergrund je nach Farbe der Einrichtung eine andere Schriftfarbe nötig ist, um ausreichenden Kontrast zu erreichen.
Durch Verwendung der Farbe \verb+AccentContrastColor+ kann hier eine automatische Anpassung im Dokument vorgenommen werden.

\subsection{Farbübersicht}

Auch die Farben der einzelnen Einheiten stehen als stabile Definitionen unabhängig von der Stiloption \verb+department+\docoption{department} zur Verfügung.

\newcommand{\farbTabellenZeile}[1]{
{\color{#1}\texttt{#1}} &
{\color{#190}\texttt{#190}} &
{\color{#150}\texttt{#150}} &
{\color{#125}\texttt{#125}} &
{\color{#1FontColor}\texttt{#1FontColor}} &
{\cellcolor{#1}\color{#1ContrastColor}\texttt{#1ContrastColor}} \\
}
\begin{center}\scriptsize
\begin{flushleft}\small\textbf{Fakultäten und Einrichtungen}\end{flushleft}\nopagebreak
\begin{tabular}{cccccc}
\farbTabellenZeile{Accent}
\farbTabellenZeile{OTHR}
\farbTabellenZeile{FakA}
\farbTabellenZeile{FakANK}
\farbTabellenZeile{FakAM}
\farbTabellenZeile{FakB}
\farbTabellenZeile{FakBM}
\farbTabellenZeile{FakBW}
\farbTabellenZeile{FakEI}
\farbTabellenZeile{FakIM}
\farbTabellenZeile{FakM}
\farbTabellenZeile{FakSG}
\farbTabellenZeile{FakS}
\farbTabellenZeile{ZWW}
\\
\farbTabellenZeile{RCAI}
\farbTabellenZeile{RCER}
\farbTabellenZeile{RCHST}
\farbTabellenZeile{RSDS}
\\
\end{tabular}

\vspace{\fill}\pagebreak

\begin{flushleft}\small\textbf{Laborfarben}\end{flushleft}\nopagebreak
\begin{tabular}{cccccc}
%A
\farbTabellenZeile{FMIS}
%ANK
\farbTabellenZeile{NACH}
\farbTabellenZeile{SappZ}
%B
\farbTabellenZeile{Geo}
\farbTabellenZeile{KNB}
%BW
\farbTabellenZeile{12Science}
\farbTabellenZeile{FEURO}
%EI
\farbTabellenZeile{Bisp}
\farbTabellenZeile{DK0PT}
\farbTabellenZeile{FENES}
\farbTabellenZeile{LAS3}
\farbTabellenZeile{MRU}
\farbTabellenZeile{SES}
%IM
\farbTabellenZeile{CCSE}
\farbTabellenZeile{eHealth}
\farbTabellenZeile{IMDS}
\farbTabellenZeile{LFD}
\farbTabellenZeile{MD}
\farbTabellenZeile{ReMIC}
\farbTabellenZeile{SEC}
%M
\farbTabellenZeile{AIMS}
\farbTabellenZeile{BFM}
\farbTabellenZeile{BMA}
\farbTabellenZeile{CEEC}
\farbTabellenZeile{CFD}
\farbTabellenZeile{FEM}
\farbTabellenZeile{IPF}
\farbTabellenZeile{KIB}
\farbTabellenZeile{KWK}
\farbTabellenZeile{LAT}
\farbTabellenZeile{LBM}
\farbTabellenZeile{LeanLab}
\farbTabellenZeile{LFT}
\farbTabellenZeile{LFW}
\farbTabellenZeile{LMP}
\farbTabellenZeile{LMS}
\farbTabellenZeile{LRT}
\farbTabellenZeile{LWS}
\farbTabellenZeile{LWM}
\farbTabellenZeile{MKS}
\farbTabellenZeile{MST}
\farbTabellenZeile{RRRU}
\farbTabellenZeile{RST}
%
\farbTabellenZeile{IST}
\farbTabellenZeile{LP}
\farbTabellenZeile{PF}
\farbTabellenZeile{PT}


\end{tabular}
\end{center}
